\section{Kiến trúc hệ thống}

\usetikzlibrary{positioning,arrows.meta,shapes.geometric,fit,calc}

\subsection{Tổng quan}

Hệ thống áp dụng mô hình kiến trúc \textit{lai} (hybrid) kết hợp Client-Server tập trung và Peer-to-Peer (P2P) phân tán.
Máy chủ trung tâm (Application Server) đảm nhiệm các chức năng xác thực, quản lý phiên, tra cứu địa chỉ ngang hàng,
quản lý cộng đồng (Community), lưu trữ và phân phối file broadcast.
Giao tiếp trực tiếp giữa hai client — bao gồm nhắn tin riêng tư và truyền file ngang hàng — được thực hiện qua kết nối TCP trực tiếp,
không qua server, nhằm giảm tải và giảm độ trễ.

\subsection{Các thành phần chính}

\begin{itemize}
    \item \textbf{Application Server}: Lắng nghe trên cổng TCP 8080 và cổng UDP 41234 (khám phá mạng).
          Mỗi kết nối TCP đến được xử lý bởi một \texttt{ClientHandler} chạy trên một \textit{virtual thread} riêng biệt.
          Server duy trì các registry trong bộ nhớ: \texttt{UserSessionRegistry} (phiên đăng nhập),
          \texttt{ChatSessionRegistry} (phiên P2P), \texttt{CommunityRegistry} (cộng đồng) và \texttt{BroadcastFileRegistry} (file broadcast).

    \item \textbf{Client}: Mỗi tiến trình client duy trì một kết nối TCP điều khiển (control socket) đến server
          và một vòng lặp đọc gói tin (\texttt{receivePackets}) chạy trên virtual thread riêng.
          Client cũng khởi động \texttt{PeerListener} — một \texttt{ServerSocket} trên cổng ngẫu nhiên — để chấp nhận kết nối P2P đến từ các peer khác.

    \item \textbf{PeerListener}: Chạy trên virtual thread, lắng nghe kết nối TCP đến từ peer.
          Phân loại kết nối dựa trên gói tin đầu tiên nhận được:
          nếu là \texttt{PEER\_HELLO} thì đây là phiên chat; nếu là \texttt{FILE\_OFFER} thì đây là luồng truyền file.

    \item \textbf{PeerConnectionRegistry}: Quản lý tập hợp các kết nối P2P đang hoạt động (\texttt{ConcurrentHashMap<String, PeerConnection>}).
          Theo dõi peer đang chọn (active peer) và ngăn chặn kết nối trùng lặp qua tập \texttt{pendingLookups}.

    \item \textbf{FileStorageService}: Lưu trữ file broadcast phía server trên hệ thống tệp cục bộ.
          Hỗ trợ upload từ client gửi và download cho các client nhận thông qua token xác thực tạm thời.

    \item \textbf{AuthService}: Quản lý xác thực người dùng. Hiện tại dữ liệu tài khoản được lưu cứng trong bộ nhớ
          (hardcoded: alice/alice123, bob/bob123, admin1/password1); chưa có cơ sở dữ liệu bền vững.
\end{itemize}

\subsection{Sơ đồ kiến trúc hệ thống}

Hình~\ref{fig:kientruc} minh họa tổng quan cách các thành phần tương tác với nhau.

\begin{figure}[H]
\centering
\begin{tikzpicture}[
    font=\small,
    box/.style={draw, rounded corners=4pt, minimum width=2.8cm, minimum height=0.75cm, align=center, fill=#1!10},
    sbox/.style={draw, rounded corners=4pt, minimum width=5.0cm, minimum height=0.7cm, align=center, fill=blue!8},
    grp/.style={draw, dashed, rounded corners=6pt, inner sep=8pt, fill=#1!5},
    arr/.style={-{Stealth[length=5pt]}, thick},
    darr/.style={{Stealth[length=5pt]}-{Stealth[length=5pt]}, thick}
]

%% ---- Client A ----
\node[box=green] (ca-con) at (-5.5, 3.2) {readConsole\\(main thread)};
\node[box=green] (ca-rcv) at (-5.5, 1.9) {receivePackets\\(virtual thread)};
\node[box=green] (ca-pl)  at (-5.5, 0.7) {PeerListener\\(virtual thread)};
\node[box=green] (ca-pr)  at (-5.5,-0.5) {PeerConnection\\Registry};

\draw[dashed, gray, rounded corners=6pt]
    (-7.3,-1.2) rectangle (-3.7, 3.85);
\node[font=\bfseries\footnotesize] at (-5.5, 4.05) {Client A};

%% ---- Server ----
\node[sbox] (s1) at (0, 3.2)  {ClientHandler (per connection)};
\node[sbox] (s2) at (0, 2.4)  {UserSessionRegistry};
\node[sbox] (s3) at (0, 1.6)  {CommunityRegistry};
\node[sbox] (s4) at (0, 0.8)  {BroadcastFileRegistry};
\node[sbox] (s5) at (0, 0.0)  {FileStorageService};
\node[sbox] (s6) at (0,-0.8)  {AuthService};

\draw[dashed, gray, rounded corners=6pt]
    (-2.8,-1.5) rectangle (2.8, 3.85);
\node[font=\bfseries\footnotesize] at (0, 4.05) {Application Server (TCP 8080 / UDP 41234)};

%% ---- Client B ----
\node[box=orange] (cb-con) at (5.5, 3.2) {readConsole\\(main thread)};
\node[box=orange] (cb-rcv) at (5.5, 1.9) {receivePackets\\(virtual thread)};
\node[box=orange] (cb-pl)  at (5.5, 0.7) {PeerListener\\(virtual thread)};
\node[box=orange] (cb-pr)  at (5.5,-0.5) {PeerConnection\\Registry};

\draw[dashed, gray, rounded corners=6pt]
    (3.7,-1.2) rectangle (7.3, 3.85);
\node[font=\bfseries\footnotesize] at (5.5, 4.05) {Client B};

%% ---- Arrows: control sockets ----
\draw[darr, blue]
    (ca-rcv.east) -- node[midway, above, font=\tiny, fill=white, inner sep=1pt] {TCP control (port 8080)}
    (s1.west |- ca-rcv);

\draw[darr, blue]
    (s1.east |- cb-rcv) -- node[midway, above, font=\tiny, fill=white, inner sep=1pt] {TCP control (port 8080)}
    (cb-rcv.west);

%% ---- Arrow: P2P ----
\draw[darr, red!70!black, dashed, thick]
    (ca-pl.east) to[out=10, in=170]
    node[midway, above=6pt, font=\tiny, fill=white, inner sep=1pt] {P2P TCP trực tiếp (chat / file)}
    (cb-pl.west);

%% ---- Broadcast upload: client A -> server ----
\draw[arr, purple]
    (-5.5,-1.0) to[out=-60, in=-140]
    node[midway, below, font=\tiny, fill=white, inner sep=1pt] {TCP mới — upload file broadcast}
    (0,-1.4);

%% ---- Broadcast download: server -> client B ----
\draw[arr, purple]
    (0,-1.4) to[out=-40, in=-120]
    node[midway, below, font=\tiny, fill=white, inner sep=1pt] {TCP mới — download file broadcast}
    (5.5,-1.0);

\end{tikzpicture}
\caption{Kiến trúc tổng quan hệ thống Client-Server kết hợp P2P}
\label{fig:kientruc}
\end{figure}

\subsection{Giao thức truyền thông}

Toàn bộ giao tiếp giữa Client và Server sử dụng giao thức frame nhị phân tùy chỉnh với cấu trúc cố định:

\[
\underbrace{[OP: 1\text{B}]}_{\text{OpCode}} \;
\underbrace{[ReqID: 4\text{B}]}_{\text{Request ID}} \;
\underbrace{[Len: 4\text{B}]}_{\text{Payload length}} \;
\underbrace{[Data: N\text{B}]}_{\text{Payload}}
\]

Phần header 9 byte được đọc trước, sau đó đọc chính xác $N$ byte payload bằng \texttt{DataInputStream.readFully()}.
Cơ chế tương quan yêu cầu-phản hồi sử dụng \texttt{ConcurrentHashMap<Integer, CompletableFuture<Packet>>} với timeout 10 giây,
cho phép nhiều yêu cầu đồng thời trên một kết nối mà không bị nhầm lẫn phản hồi.

\subsection{Quản lý luồng thực thi}

Hệ thống sử dụng \textit{virtual threads} (Project Loom, Java 21+) thay vì platform threads truyền thống để xử lý đồng thời:

\begin{itemize}
    \item Mỗi kết nối TCP đến server tạo một virtual thread chạy \texttt{ClientHandler}.
    \item Mỗi \texttt{PeerConnection} có một virtual thread đọc gói tin độc lập.
    \item \texttt{PeerListener} chạy trên một virtual thread riêng; mỗi kết nối mới được phân phối ngay cho virtual thread tiếp theo.
    \item Trên client, \texttt{receivePackets} và \texttt{PeerListener} đều là virtual thread, song song với luồng chính đọc console.
\end{itemize}

Việc đồng bộ hóa ghi lên socket được thực hiện qua khóa đối tượng (object lock): \texttt{serverWriteLock} phía client,
\texttt{synchronized(this)} trong \texttt{ClientHandler.sendPacket()}, và \texttt{synchronized(outputStream)} trong vòng lặp broadcast phía server.

\subsection{Quản lý cộng đồng}

Cộng đồng (Community) là khái niệm nhóm người dùng do server quản lý tập trung.
\texttt{CommunityRegistry} lưu trữ các đối tượng \texttt{Community} trong \texttt{ConcurrentHashMap} theo ID.
Mỗi Community có danh sách thành viên; chỉ thành viên mới nhận được tin nhắn broadcast và file broadcast của cộng đồng đó.
Mọi trạng thái cộng đồng hiện tại chỉ tồn tại trong bộ nhớ và bị mất khi server khởi động lại.
